\documentclass[11pt]{amsart}

\usepackage[T1]{fontenc}
\usepackage{lmodern}
\usepackage{microtype}
\usepackage{mathtools,amssymb,amsthm}
\usepackage[hidelinks]{hyperref}

\hypersetup{
  pdftitle={The Jaccard Triangle Inequality for Supermodular Log-Submodular
            Valuations on Arbitrary Lattices}
}

\newtheorem{theorem}{Theorem}
\newtheorem{lemma}[theorem]{Lemma}
\newtheorem{corollary}[theorem]{Corollary}
\newtheorem{proposition}[theorem]{Proposition}
\theoremstyle{remark}
\newtheorem{remark}[theorem]{Remark}

\newcommand{\Lat}{L}
\newcommand{\dJ}{d_{f,J}}

\title[Jaccard triangle inequality on arbitrary lattices]
{The Jaccard Triangle Inequality for Supermodular Log-Submodular Valuations
on Arbitrary Lattices}

\date{}

\subjclass[2020]{06B05, 06B99, 54E35}
\keywords{generalized Jaccard distance, Tanimoto distance, supermodular
valuation, log-submodular valuation, lattice, triangle inequality,
non-distributive lattice}

\begin{document}

\begin{abstract}
Let $\Lat$ be a lattice and $f\colon \Lat\to(0,\infty)$ a valuation. The
generalized Jaccard distance is
$\dJ(A,B)=1-f(A\wedge B)/f(A\vee B)$.
Badica and Badica prove that $\dJ$ satisfies the triangle inequality when $f$
is strictly positive, monotone, supermodular and log-submodular, provided
$\Lat$ is relatively complemented and distributive, and ask in their
Section~7 whether those two structural hypotheses on $\Lat$ can be removed. We
observe that they can: the triangle inequality holds on an arbitrary lattice,
with no distributivity, relative complementation, bounds or finiteness assumed.
The proof applies supermodularity once and log-submodularity once, not to
$A,B,C$ but to the derived pairs $(A\wedge B,\,B\wedge C)$ and
$(A\vee B,\,B\vee C)$, and collapses to the inequality
$(P-p)(Q-q)\ge(b-p)(b-q)$ between five values of $f$. Dropping either
supermodularity or log-submodularity breaks the conclusion already on the
four-element Boolean lattice.
\end{abstract}

\maketitle

\section{The question}

Let $\Lat$ be a lattice, written with $\wedge$ and $\vee$ and ordered by
$X\le Y\iff X\wedge Y=X$. Call $f\colon\Lat\to(0,\infty)$
\emph{monotone} if $X\le Y\Rightarrow f(X)\le f(Y)$,
\emph{supermodular} if
\begin{equation}\label{eq:super}
 f(X)+f(Y)\ \le\ f(X\vee Y)+f(X\wedge Y)
 \qquad(X,Y\in \Lat),
\end{equation}
and \emph{log-submodular} if
\begin{equation}\label{eq:logsub}
 f(X)\,f(Y)\ \ge\ f(X\vee Y)\,f(X\wedge Y)
 \qquad(X,Y\in \Lat).
\end{equation}
The \emph{generalized Jaccard distance} of $f$ is
\begin{equation}\label{eq:jac}
 \dJ(A,B)\ =\ 1-\frac{f(A\wedge B)}{f(A\vee B)},
 \qquad
 S(A,B)\ :=\ \frac{f(A\wedge B)}{f(A\vee B)},
\end{equation}
so $\dJ=1-S$, and the triangle inequality
$\dJ(A,C)\le\dJ(A,B)+\dJ(B,C)$ is exactly
\begin{equation}\label{eq:TI}
 S(A,B)+S(B,C)\ \le\ 1+S(A,C).
\end{equation}

In the preprint \cite{BB}, Badica and Badica prove \eqref{eq:TI} for strictly
positive monotone supermodular log-submodular $f$ on a lattice that is
relatively complemented \emph{and} distributive (their supermodular
log-submodular theorem), and separately prove \eqref{eq:TI} for modular $f$ on
an arbitrary lattice (their modular theorem). Their Section~7, \emph{Discussion
of Open Theoretical Problems}, lists four open problems; the second of them
reads verbatim and in full:

\begin{quote}\sloppy
\textbf{Arbitrary lattices for supermodular valuations:} While we established
metricity for supermodular and $\log$-submodular valuations, the proof relies
on the lattice being both relatively complemented and distributive. It is
still unknown whether the standard generalized Jaccard formula holds up to the
triangle inequality for these valuations on arbitrary lattices that lack those
two specific traits.
\end{quote}

Here ``the standard generalized Jaccard formula'' is \eqref{eq:jac} and
``these valuations'' is the hypothesis class of their supermodular
log-submodular theorem, namely strictly positive, monotone, supermodular and
log-submodular. No conjecture is attached to the bullet, so what follows
answers it in the affirmative and neither confirms nor refutes a stated guess.

\section{The theorem}

\begin{theorem}\label{thm:main}
Let $\Lat$ be a lattice, not assumed distributive, relatively complemented,
bounded or finite, and let $f\colon\Lat\to(0,\infty)$ be monotone,
supermodular and log-submodular. Then $\dJ$ satisfies the triangle inequality: for all
$A,B,C\in \Lat$,
\[
 \dJ(A,C)\ \le\ \dJ(A,B)+\dJ(B,C).
\]
\end{theorem}

Hence the supermodular log-submodular theorem of \cite{BB} holds verbatim with
the words ``relatively complemented distributive'' deleted, which answers the
bullet quoted above.

\begin{proof}
Fix $A,B,C\in \Lat$ and abbreviate the seven values of $f$ involved:
\[
 p=f(A\wedge B),\quad P=f(A\vee B),\quad
 q=f(B\wedge C),\quad Q=f(B\vee C),
\]
\[
 r=f(A\wedge C),\quad R=f(A\vee C),\quad
 b=f(B).
\]
All seven are $>0$. Since $A\wedge B\le B\le A\vee B$ and
$B\wedge C\le B\le B\vee C$, monotonicity gives
\begin{equation}\label{eq:mono}
 0<p\le b\le P,
 \qquad
 0<q\le b\le Q .
\end{equation}
By \eqref{eq:TI} it suffices to prove $p/P+q/Q\le 1+r/R$.

\smallskip
\noindent\textbf{Step 1 (supermodularity, on the pair of meets).}
Put $X=A\wedge B$ and $Y=B\wedge C$. Both lie below $B$, so $X\vee Y\le B$;
and $X\wedge Y=A\wedge B\wedge C\le A\wedge C$. Applying \eqref{eq:super} to
this pair and then monotonicity twice,
\[
 p+q\ =\ f(X)+f(Y)\ \le\ f(X\vee Y)+f(X\wedge Y)\ \le\ b+r,
\]
that is
\begin{equation}\label{eq:lem1}
 r\ \ge\ p+q-b .
\end{equation}

\smallskip
\noindent\textbf{Step 2 (log-submodularity, on the pair of joins).}
Put $X=A\vee B$ and $Y=B\vee C$. Both lie above $B$, so $X\wedge Y\ge B$; and
$X\vee Y=A\vee B\vee C\ge A\vee C$. Applying \eqref{eq:logsub} to this pair
and then monotonicity twice,
\[
 PQ\ =\ f(X)f(Y)\ \ge\ f(X\vee Y)\,f(X\wedge Y)\ \ge\ R\,b,
\]
that is
\begin{equation}\label{eq:lem2}
 R\ \le\ \frac{PQ}{b}.
\end{equation}

\smallskip
\noindent\textbf{Step 3 (a lower bound for $S(A,C)$).}
We claim $r/R\ \ge\ b(p+q-b)/(PQ)$. If $p+q-b\le 0$ the right-hand side is
$\le0<r/R$. If $p+q-b>0$ then $r\ge p+q-b>0$ by \eqref{eq:lem1} and
$0<R\le PQ/b$ by \eqref{eq:lem2}, so
$r/R\ge (p+q-b)\,b/(PQ)$. In either case it suffices to prove
\begin{equation}\label{eq:suff}
 \frac{p}{P}+\frac{q}{Q}\ \le\ 1+\frac{b\,(p+q-b)}{PQ}.
\end{equation}

\smallskip
\noindent\textbf{Step 4 (clearing denominators).}
Multiplying \eqref{eq:suff} by $PQ>0$ turns it into
$pQ+qP\le PQ+bp+bq-b^2$, and
\[
 \bigl(pQ+qP\bigr)-\bigl(PQ+bp+bq-b^{2}\bigr)
 \ =\ -\Bigl[(P-p)(Q-q)-(b-p)(b-q)\Bigr]
\]
identically in $p,P,q,Q,b$ --- both sides expand to
$pQ+Pq-PQ-bp-bq+b^{2}$. So \eqref{eq:suff} is precisely
\begin{equation}\label{eq:prod}
 (P-p)(Q-q)\ \ge\ (b-p)(b-q).
\end{equation}

\smallskip
\noindent\textbf{Step 5 (a termwise certificate for \eqref{eq:prod}).}
Identically in $p,P,q,Q,b$,
\begin{equation}\label{eq:cert}
 (P-p)(Q-q)-(b-p)(b-q)\ =\ (P-b)(Q-q)+(b-p)(Q-b).
\end{equation}
By \eqref{eq:mono} each of $P-b$, $Q-q$, $b-p$, $Q-b$ is $\ge0$, so the
right-hand side of \eqref{eq:cert} is a sum of two products of non-negative
numbers and is $\ge0$. This proves \eqref{eq:prod}, hence \eqref{eq:suff},
hence \eqref{eq:TI}.
\end{proof}

\begin{remark}\sloppy
The proof uses only the four inequalities \eqref{eq:mono} and one application
each of \eqref{eq:super} and \eqref{eq:logsub}, to the derived pairs
$(A\wedge B,\,B\wedge C)$ and $(A\vee B,\,B\vee C)$ rather than to $(A,B)$,
$(B,C)$ or $(A,C)$. The proof in \cite{BB} instead constructs relative
complements inside the interval $[A\wedge B\wedge C,\ A\vee B\vee C]$ and then
applies distributivity, which is where its two structural hypotheses on $\Lat$
enter.
\end{remark}

\begin{remark}\label{rem:pseudo}
Theorem~\ref{thm:main} is about the triangle inequality, which is what the
bullet asks about. It does not upgrade $\dJ$ to a metric: $\dJ(X,Y)=0$ exactly
when $f$ is constant on $[X\wedge Y,\,X\vee Y]$, so $\dJ$ is a pseudometric in
general and a metric exactly when $f$ is strictly monotone. The same caveat
attaches to the theorems of \cite{BB}.
\end{remark}

\section{Both hypotheses on $f$ are needed}

Write $B_2$ for the four-element Boolean lattice, given by its covers:
\begin{equation}\label{eq:B2}
 B_2:\quad 0<x,\ \ 0<y,\ \ x<1,\ \ y<1 .
\end{equation}
Write a valuation on $B_2$ as the tuple $\bigl(f(0),f(x),f(y),f(1)\bigr)$, and
recall $\Delta(A,B,C):=S(A,B)+S(B,C)-1-S(A,C)$, so that $\Delta>0$ is a
violation of \eqref{eq:TI}. Each row below is an exact rational computation on
the object printed in \eqref{eq:B2}.
\begin{center}
\begin{tabular}{lcccl}
$f$ & lost & $(A,B,C)$ & $\Delta$ & step that breaks\\
\hline
$(2,3,3,20)$ & \eqref{eq:logsub} & $(x,0,y)$ & $\tfrac{7}{30}$ & 2:
 $Rb=40>9=PQ$\\
$(1,3,3,4)$  & \eqref{eq:super}  & $(x,1,y)$ & $\tfrac14$ & 1:
 $p+q-b=2>1=r$
\end{tabular}
\end{center}
The first valuation is monotone and supermodular but not log-submodular
($f(1)f(0)=40>9=f(x)f(y)$); the second is monotone and log-submodular but not
supermodular ($f(1)+f(0)=5<6=f(x)+f(y)$). So neither supermodularity nor
log-submodularity may be dropped from Theorem~\ref{thm:main}, and neither is
implied by monotonicity; and since $B_2$ is distributive and relatively
complemented, this bounds the hypotheses on $f$ and not the hypotheses on
$\Lat$ that Theorem~\ref{thm:main} removes.

Theorem~\ref{thm:main} is a sufficiency statement. Section~6 of \cite{BB}
proves, on an arbitrary lattice, that monotonicity together with
supermodularity of $f$ \emph{and} of $1/f$ is \emph{necessary} for
\eqref{eq:TI}; that is a different hypothesis class, so
Theorem~\ref{thm:main} is neither the converse of it nor half of a
characterisation. The remaining open problems of Section~7 of \cite{BB} are
untouched here.

\section{Verification}\label{sec:verify}

The proof of Theorem~\ref{thm:main} is a hand argument and needs no machine.
The accompanying program \texttt{verify.py} (Python 3.9 or later, standard
library only; exact arithmetic throughout, no floating point) re-derives the
numbers stated here: it verifies the two polynomial identities of Steps~4
and~5 as exact identities in $\mathbb{Z}[p,P,q,Q,b]$ by sparse expansion
rather than by sampling; it confirms the four order facts used in Steps~1
and~2 at every ordered triple of every lattice of order at most six; it checks
that the covers \eqref{eq:B2} define a lattice; and it reproduces the two
violation margins $7/30$ and $1/4$ tabulated above together with the step each
breaks. It prints one line per check and closes with
\texttt{VERDICT: ALL 27 CHECKS PASS}, exiting $0$ only if all pass.

Most of the 27 checks bear no load on Theorem~\ref{thm:main}: finite sweeps
over small lattices, an enumerator control and further example valuations are
recorded but are not used above, and no program can quantify over the class the
theorem covers. The recorded output is also not labelled for the present text:
it letters its own sections A--G, calls the inequalities \eqref{eq:lem1}
$r\ge p+q-b$ and \eqref{eq:lem2} $R\le PQ/b$ of Steps~1 and~2 \emph{Lemma 1}
and \emph{Lemma 2}, and describes three lattices, three totals and a valuation
$g$ as printed in Sections~3 and~4 of the paper. None of those cross-references
holds here: the only objects this paper prints are the lattice \eqref{eq:B2}
and the two valuations tabulated in Section~3, it states no totals, and the
valuation $g$ appears only in the program. The theorems of \cite{BB} are quoted
there, not re-proved.

\begin{thebibliography}{1}

\bibitem{BB}
C.~Badica and A.~Badica,
\emph{On the triangle inequality for the Jaccard distance in arbitrary
lattices},
arXiv:2608.18194v1, 2026.
\href{https://arxiv.org/abs/2608.18194}{arXiv:2608.18194}.
The passages quoted above are from the v1 e-print, the only version available
at the time of writing.

\end{thebibliography}

\end{document}
